\documentclass[11pt,letterpaper]{article}
% ==============================================================================
% Shared Preamble: Principles of Economics Modular Curriculum
% Core Colors: SlateBlue (#1E293B) & Terracotta (#C2410C)
% ==============================================================================

\usepackage[utf8]{inputenc}
\usepackage[T1]{fontenc}
\usepackage{lmodern}
\usepackage[margin=1in]{geometry}
\usepackage{amsmath,amssymb,amsfonts,amsthm}
\usepackage{booktabs}
\usepackage{microtype}
\usepackage{enumitem}
\usepackage{fancyhdr}
\usepackage{titlesec}
\usepackage{xcolor}
\usepackage{tikz}
\usepackage{pgfplots}
\usepackage{mdframed}
\usepackage[hidelinks,colorlinks=true,linkcolor=SlateBlue,citecolor=SlateBlue,urlcolor=Terracotta]{hyperref}

\pgfplotsset{compat=1.18}
\usetikzlibrary{arrows.meta,positioning,calc,patterns,decorations.pathreplacing}

% --- Color Palette Definitions ---
\definecolor{SlateBlue}{HTML}{1E293B}      % Primary Accent / Structure
\definecolor{Terracotta}{HTML}{C2410C}     % Secondary Accent / Emphasis
\definecolor{ForestGreen}{HTML}{15803D}    % Equilibrium / Positive Shocks
\definecolor{LightSlate}{HTML}{F8FAFC}     % Box Background Fill
\definecolor{BorderGray}{HTML}{CBD5E1}     % Neutral Borders
\definecolor{LightTerracotta}{HTML}{FFF7ED}% Cautionary / Principle Fill

% --- Section Heading Styling ---
\titleformat{\section}
  {\normalfont\Large\bfseries\color{SlateBlue}}{\thesection}{1em}{}[\titlerule]
\titleformat{\subsection}
  {\normalfont\large\bfseries\color{SlateBlue}}{\thesubsection}{1em}{}
\titleformat{\subsubsection}
  {\normalfont\normalsize\bfseries\color{Terracotta}}{\thesubsubsection}{1em}{}

% --- Header and Footer Configuration ---
\pagestyle{fancy}
\fancyhf{}
\renewcommand{\headrulewidth}{0.4pt}
\renewcommand{\footrulewidth}{0.4pt}
\fancyhead[L]{\small\scshape\color{SlateBlue}\leftmark}
\fancyhead[R]{\small\color{SlateBlue}Principles of Economics}
\fancyfoot[C]{\thepage}

% --- Custom Environments (mdframed) ---

% 1. Core Principle Box
\newmdenv[
  linecolor=Terracotta,
  linewidth=1.5pt,
  backgroundcolor=LightTerracotta,
  innertopmargin=10pt,
  innerbottommargin=10pt,
  innerrightmargin=12pt,
  innerleftmargin=12pt,
  roundcorner=4pt,
  skipabove=12pt,
  skipbelow=12pt
]{coreprinciple}

% 2. Empirical Case Study Box
\newmdenv[
  linecolor=SlateBlue,
  linewidth=1.5pt,
  backgroundcolor=LightSlate,
  innertopmargin=10pt,
  innerbottommargin=10pt,
  innerrightmargin=12pt,
  innerleftmargin=12pt,
  roundcorner=4pt,
  skipabove=12pt,
  skipbelow=12pt
]{casestudy}

% 3. Key Takeaways Box
\newmdenv[
  linecolor=BorderGray,
  linewidth=1pt,
  backgroundcolor=white,
  innertopmargin=10pt,
  innerbottommargin=10pt,
  innerrightmargin=12pt,
  innerleftmargin=12pt,
  roundcorner=4pt,
  skipabove=12pt,
  skipbelow=12pt
]{takeaways}

% Environment Title Helpers
\newcommand{\boxheader}[2]{{\noindent\bfseries\color{#1}\Large #2}\vspace{0.4em}\par}
\usepackage{tikz}
\usepackage{pgfplots}
\pgfplotsset{compat=1.18}
\usetikzlibrary{arrows.meta,positioning,calc,shapes.geometric,patterns,decorations.pathreplacing,intersections}

\tikzset{
  econ axis/.style={->, >=Stealth, thick, color=SlateBlue!90!black},
  econ grid/.style={thin, dashed, color=gray!30},
  econ guide/.style={dashed, color=gray!60, thin},
  demand curve/.style={very thick, color=SlateBlue},
  supply curve/.style={very thick, color=Terracotta},
  mb curve/.style={very thick, color=SlateBlue},
  mc curve/.style={very thick, color=Terracotta},
  econ point/.style={circle, fill=SlateBlue!90!black, inner sep=1.8pt},
  econ highlight/.style={fill=Terracotta!20, opacity=0.5},
  econ label/.style={font=\small\sffamily, align=center},
  econ annotation/.style={font=\footnotesize\itshape, color=SlateBlue!80!black}
}


\begin{document}

\title{\Huge\bfseries\color{SlateBlue} Chapter 6: Saving, Investment Spending, and the Financial System}
\author{\large Principles of Macroeconomics Curriculum}
\date{2026}
\maketitle

\thispagestyle{plain}

\begin{abstract}
This chapter provides a rigorous, pedagogical treatment of Chapter 6: Saving, Investment Spending, and the Financial System, benchmarked to standard principles curriculum.
\end{abstract}

\tableofcontents
\vspace{1.5em}
\hrule
\vspace{1.5em}

\section{Learning Outcomes and Key Terms}
Students should distinguish saving from wealth, apply the national saving identity, and use the loanable-funds model to explain how saving and investment plans respond to real interest rates. These objectives support MA-2 and MA-3 and prepare students for C2002-2. Key terms include private saving, public saving, national saving, investment, loanable funds, real interest rate, deficit, and capital inflow.

\section{Saving Is a Flow; Wealth Is a Stock}
A household that saves part of this month's income adds to its assets, but saving itself is measured over a period. Wealth is the value of assets held at a point in time. Likewise, investment in macroeconomics is spending on newly produced capital goods and inventories during a period; it is not the purchase of an existing stock certificate. Keeping stocks and flows separate prevents a common accounting error.

\section{National Saving and the Open-Economy Identity}
Start with the expenditure identity $Y=C+I+G+NX$. Rearranging gives $Y-C-G=I+NX$. National saving is $S=Y-C-G$, so in the simplified model $S=I+NX$. If net exports are used as the current-account proxy, then $S-I=NX$: saving above domestic investment corresponds to net lending abroad, while investment above saving is matched by net borrowing from abroad. In full balance-of-payments accounts, the current account includes net primary income and transfers as well as goods and services.

With taxes $T$ and disposable income $Y-T$, private saving is $S_p=Y-T-C$ and public saving is $S_g=T-G$. National saving adds the two:
\begin{equation}
S=S_p+S_g=(Y-T-C)+(T-G)=Y-C-G, \qquad S-I=NX \quad\text{(simplified model)}.
\end{equation}
In a closed economy, $NX=0$ and planned national saving equals investment in the accounting identity. The identity does not guarantee that desired saving and planned investment are equal at every interest rate; prices, output, and financial markets adjust toward consistency.

\section{Loanable Funds: Graph and Mechanism}
Put the quantity of loanable funds on the horizontal axis and the real interest rate on the vertical axis. Planned saving is drawn as a supply of funds, often rising with the real return. Planned investment is a demand for funds, often falling as borrowing becomes more expensive. Their intersection is the model's equilibrium real rate and quantity. A shift in saving or investment plans moves one curve; a change in the interest rate is a movement along a curve. This is a market model, not a literal inventory of all bank deposits.

\section{Deficits, Interest Rates, and Crowding Out}
If government purchases rise while taxes and private saving behavior are fixed, public saving falls. In a closed-economy loanable-funds model, the saving supply can shift left, raising the equilibrium real interest rate and reducing some private investment. This is the crowding-out mechanism. Its size depends on conditions: monetary policy, spare capacity, household saving responses, international capital flows, and the central bank's reaction can change the outcome. A deficit does not mechanically imply a one-for-one decline in private investment.

\begin{examplecase}{FRED Data: Gross Private Domestic Investment}
FRED series GPDI reports gross private domestic investment in billions of current dollars at a seasonally adjusted annual rate, quarterly. It includes fixed investment and inventory change. It is a nominal expenditure flow, not a capital stock and not a direct measure of productive capacity. To study physical investment growth across time, compare an appropriate real series and account for revisions and price changes.
\end{examplecase}

\begin{coreprinciple}
\boxheader{SlateBlue}{Clear It Up: The Identity Is Not a Causal Claim}
The accounting relation $S-I=NX$ in the simplified model must hold ex post. It does not by itself tell us whether saving caused investment, investment caused income and saving, or foreign flows adjusted. Use a behavioral model and evidence to discuss causation.
\end{coreprinciple}

\section{Check Your Understanding}
\begin{enumerate}
  \item If $Y=1{,}000$, $C=650$, and $G=200$, calculate national saving.
  \item If national saving is 150 and investment is 180, what is $NX$ in the simplified open-economy identity?
  \item Is buying an existing share of a company macroeconomic investment? Why or why not?
\end{enumerate}
Answers: (1) $S=Y-C-G=150$. (2) $NX=S-I=-30$, a deficit in the simplified model. (3) No; it transfers ownership of an existing financial asset rather than purchasing newly produced capital.

\end{document}
